\chapter{Background and Related Work}
\label{ch:background}

This chapter provides the technical background necessary to understand the design
and evaluation of Ai4k8s. Section~\ref{sec:kubernetes} introduces Kubernetes and
its autoscaling primitives. Section~\ref{sec:aiops} surveys the AIOps paradigm.
Section~\ref{sec:llm-reasoning} examines the use of \glspl{LLM} as reasoning
engines. Section~\ref{sec:mcda} introduces multi-criteria decision analysis.
Section~\ref{sec:related} positions Ai4k8s in the related work.

% -----------------------------------------------------------------------
\section{Kubernetes and Container Orchestration}
\label{sec:kubernetes}

Kubernetes~\cite{k8s2014} is an open-source container orchestration platform
originally developed at Google and donated to the Cloud Native Computing Foundation
in 2016. It automates the deployment, scaling, and lifecycle management of
containerized applications organized into \textit{pods}—the smallest deployable
unit, comprising one or more containers sharing a network namespace and storage.

A Kubernetes cluster consists of a \textit{control plane} (API server, scheduler,
controller manager, etcd) and one or more \textit{worker nodes} that execute pods.
Resource allocation is managed through \textit{requests} and \textit{limits}
specified per container: requests determine scheduling placement, while limits
cap consumption.

\subsection{Horizontal Pod Autoscaler (HPA)}
\label{sec:hpa}

The \gls{HPA}~\cite{k8s-hpa} adjusts the number of pod replicas in a Deployment
or ReplicaSet based on observed resource metrics. The default control loop runs
every 15~seconds and computes a target replica count as:

\begin{equation}
  \text{desiredReplicas} = \left\lceil
    \text{currentReplicas} \times
    \frac{\text{currentMetricValue}}{\text{desiredMetricValue}}
  \right\rceil
  \label{eq:hpa}
\end{equation}

Scaling is subject to configurable \texttt{minReplicas} and \texttt{maxReplicas}
bounds and a stabilization window (default 300~s for scale-down) that prevents
thrashing. The HPA is purely reactive: it responds only to metrics that have
already exceeded or fallen below the configured threshold. Any provisioning lag
between metric breach and pod readiness is invisible to the controller.

\subsection{Vertical Pod Autoscaler (VPA)}
\label{sec:vpa}

The \gls{VPA}~\cite{k8s-vpa} right-sizes container resource requests and limits
based on observed usage history rather than adding replicas. It operates in three
modes:
\begin{itemize}
  \item \texttt{Recommender}: computes resource recommendations from historical
  percentile usage without applying them.
  \item \texttt{Updater}: evicts pods whose current requests deviate significantly
  from the recommendation (requires pod restart).
  \item \texttt{Admission Controller}: patches new pod resource specs at admission
  time with the latest recommendation.
\end{itemize}
VPA and HPA are not recommended to be used together on the same resource metric,
as their control loops can conflict. VPA is best suited to stateful workloads where
horizontal scaling is impractical. Its key limitation is the observation window:
a fresh deployment has no history, and the Recommender typically requires several
minutes of load before producing its first recommendation.

\subsection{Limitations of Reactive Autoscaling}
\label{sec:reactive-limits}

Both HPA and VPA share structural limitations that motivate a predictive approach:

\begin{enumerate}
  \item \textbf{Provisioning lag.} Pod startup time (image pull, init containers,
  readiness probe) ranges from seconds to minutes. A reactive autoscaler that
  triggers at threshold breach may not have healthy replicas available until the
  peak has passed.

  \item \textbf{Single-criterion optimization.} HPA optimizes a single metric;
  the operational reality of balancing cost, latency, and stability requires
  multi-dimensional reasoning.

  \item \textbf{No forecasting.} Neither mechanism uses predicted future load.
  Seasonal patterns, scheduled events, and known traffic ramps are ignored.

  \item \textbf{No explainability.} Neither mechanism records why a scaling
  decision was made. Post-hoc auditing requires reconstructing metric history.
\end{enumerate}

% -----------------------------------------------------------------------
\section{AIOps: AI for IT Operations}
\label{sec:aiops}

\gls{AIOps} is defined by Gartner~\cite{GartnerAIOpsDefinition} as the application
of machine learning and big data analytics to IT operations to augment or replace
processes driven by human expertise. AIOps systems typically ingest telemetry
streams (metrics, logs, traces), apply analytical models to detect anomalies,
predict failures, correlate events, and automate remediation.

Early AIOps systems focused on anomaly detection using statistical methods
(ARIMA, Isolation Forest) and supervised classifiers trained on labelled incident
data~\cite{zhang2019robust}. More recent work has moved toward deep learning
for log parsing~\cite{du2017deeplog} and root cause analysis. The introduction
of \glspl{LLM} has opened a new frontier: models with broad operational knowledge
that can reason across heterogeneous signals and produce natural-language
explanations~\cite{cheng2023ai}.

Remaining challenges include data quality and label scarcity in production
environments, the computational cost of running large models continuously, and
the trust deficit arising from opaque model outputs—precisely the gap that
explainability mechanisms aim to close.

% -----------------------------------------------------------------------
\section{Large Language Models as Reasoning Engines}
\label{sec:llm-reasoning}

\glspl{LLM} are transformer-based neural networks trained on large corpora of
text to predict the next token in a sequence~\cite{brown2020gpt3}. Models such
as GPT-4, LLaMA, and the Qwen family have demonstrated emergent capabilities
including in-context learning, chain-of-thought reasoning, and structured output
generation—properties that make them attractive as reasoning engines in agentic
AI systems.

\subsection{Structured Output and Function Calling}
For operational decision-making, an LLM must produce reliably parseable output.
Modern APIs support \textit{structured output} modes that constrain the model
to emit valid JSON matching a specified schema. This eliminates the need for
fragile regex parsing of free-text responses and enables downstream components
to consume LLM decisions programmatically.

\subsection{Local Inference with Quantized Models}
Running LLMs on-premises avoids data privacy concerns, eliminates API latency
variance, and removes cloud dependency. Quantization reduces model size by
lowering the numeric precision of weights (e.g., from FP16 to INT4), enabling
inference on commodity hardware without a GPU. The GGUF format~\cite{ggml2023}
is widely used for CPU inference. Table~\ref{tab:model-sizes} shows the models
evaluated in this work.

\begin{table}[h]
  \centering
  \caption{Qwen model variants evaluated for on-premises autoscaling inference.}
  \label{tab:model-sizes}
  \begin{tabular}{llrrl}
    \toprule
    Model & Quantization & Disk & RAM (loaded) & Idle inference \\
    \midrule
    qwen3.5:2b  & Q8\_0 & 2.7~GB & $\sim$3.8~GiB & 7.7~s \\
    qwen3.5:2b  & Q4\_K\_M & 1.9~GB & $\sim$2.3~GiB & 11~s \\
    qwen3:0.6b  & Q4   & 522~MB & $\sim$0.7~GiB & \textbf{1.7~s} \\
    \bottomrule
  \end{tabular}
\end{table}

\subsection{LLM Cascades}
A cascade architecture chains multiple inference backends in order of preference,
falling back to the next when the primary is unavailable or too slow. This pattern
is used in Ai4k8s to combine local (Ollama/Qwen) and cloud (Groq) inference:
local inference runs at idle; the cascade activates the cloud backend when local
inference exceeds a configurable soft timeout.

% -----------------------------------------------------------------------
\section{Multi-Criteria Decision Analysis (MCDA)}
\label{sec:mcda}

\gls{MCDA} encompasses a family of methods for evaluating alternatives across
multiple conflicting criteria. In resource allocation problems, typical criteria
include cost, performance, reliability, and risk~\cite{hwang1981topsis}.

\subsection{TOPSIS}
\label{sec:topsis}

The Technique for Order of Preference by Similarity to Ideal Solution
(TOPSIS)~\cite{hwang1981topsis} selects the alternative closest to the positive
ideal solution (best possible value on every criterion) and farthest from the
negative ideal solution (worst possible value on every criterion).

Let $\mathbf{X} \in \mathbb{R}^{m \times n}$ denote the decision matrix with $m$
alternatives and $n$ criteria. TOPSIS proceeds as follows:

\begin{enumerate}
  \item \textbf{Normalization.} Each entry is normalized to remove scale differences:
  \begin{equation}
    r_{ij} = \frac{x_{ij}}{\sqrt{\sum_{i=1}^{m} x_{ij}^2}}.
    \label{eq:topsis-norm}
  \end{equation}

  \item \textbf{Weighting.} Normalized values are weighted by criterion importance:
  \begin{equation}
    v_{ij} = w_j \cdot r_{ij}, \qquad \sum_{j=1}^{n} w_j = 1.
    \label{eq:topsis-weight}
  \end{equation}

  \item \textbf{Ideal and anti-ideal solutions.} For each criterion $j$:
  \begin{equation}
    A^{+}_j =
    \begin{cases}
      \max_i v_{ij}, & \text{benefit criterion,} \\
      \min_i v_{ij}, & \text{cost criterion,}
    \end{cases}
    \qquad
    A^{-}_j \text{ defined conversely.}
    \label{eq:topsis-ideal}
  \end{equation}

  \item \textbf{Separation measures.} Euclidean distances to ideal and anti-ideal:
  \begin{equation}
    D^{+}_i = \lVert \mathbf{v}_i - A^{+} \rVert_2, \qquad
    D^{-}_i = \lVert \mathbf{v}_i - A^{-} \rVert_2.
    \label{eq:topsis-dist}
  \end{equation}

  \item \textbf{Closeness coefficient.} The relative closeness to the ideal:
  \begin{equation}
    C_i = \frac{D^{-}_i}{D^{+}_i + D^{-}_i} \in [0,1].
    \label{eq:topsis-closeness}
  \end{equation}
  The alternative with the highest $C_i$ is selected.
\end{enumerate}

TOPSIS is computationally lightweight (linear in the number of alternatives and
criteria), deterministic, and produces an interpretable score for each alternative
that can be logged as part of a decision trace.

\paragraph{MCDA Override Mechanism.}
In Ai4k8s, the MCDA module compares the LLM-recommended target $r_L$ with the
TOPSIS-optimal target $r_M$. If the TOPSIS closeness coefficient of $r_M$ exceeds
that of $r_L$ by more than a dominance threshold $\delta = 0.15$, the MCDA
recommendation overrides the LLM. The agreement is classified as:
\begin{itemize}
  \item \textbf{Full agreement:} same scaling direction and $|r_L - r_M| \leq 1$.
  \item \textbf{Partial agreement:} same direction, different magnitude.
  \item \textbf{Divergence/override:} opposite directions or gap $> \delta$.
\end{itemize}

% -----------------------------------------------------------------------
\section{Time-Series Forecasting}
\label{sec:forecasting}

Ai4k8s uses time-series forecasting to inject a look-ahead signal into the LLM
context. Two forecast models are evaluated: \gls{ARIMA} and \gls{LSTM}.

\textbf{ARIMA}~\cite{box1970arima} combines autoregressive (AR), integrating (I),
and moving-average (MA) components. It is effective for stationary or
difference-stationary series with a moderate number of lags and is computationally
cheap to fit incrementally.

\textbf{LSTM}~\cite{hochreiter1997lstm} is a recurrent neural network architecture
with gating mechanisms that address the vanishing gradient problem. LSTMs capture
long-range temporal dependencies and perform well on nonlinear, multi-pattern
workloads, at the cost of higher training complexity.

Ai4k8s's predictive monitoring system maintains a 288-sample (24-hour) ring
buffer and selects the lower-MAPE model dynamically per deployment.

% -----------------------------------------------------------------------
\section{Related Work}
\label{sec:related}

\textbf{Kubernetes autoscaling surveys.} Ilyushkin~et~al.~\cite{ilyushkin2017experimental}
compare threshold-based and model-predictive autoscalers on cloud workloads,
finding that predictive controllers reduce SLA violations under periodic load
patterns. Rossi~et~al.~\cite{rossi2019horizontal} analyze HPA stability margins
and propose hysteresis-based improvements. Neither work incorporates LLM reasoning
or multi-criteria optimization.

\textbf{LLM-based infrastructure management.} AIOpsLab~\cite{chen2024aiopslab}
provides a benchmark for evaluating LLM agents on incident response and root cause
analysis tasks. K8sGPT~\cite{k8sgpt2023} uses LLMs to explain Kubernetes diagnostic
output. Neither system makes autonomous scaling decisions. ReactAgent~\cite{yao2023react}
introduces the ReAct pattern (Reasoning + Acting), showing that LLMs interleaved
with tool calls outperform pure chain-of-thought on planning tasks—a conceptual
ancestor of the Ai4k8s reasoning layer.

\textbf{MCDA in cloud resource management.} Taghavi~et~al.~\cite{taghavi2010resource}
apply TOPSIS to virtual machine placement in cloud data centers, demonstrating
significant cost savings over greedy first-fit scheduling. Our work extends this
to dynamic autoscaling within a running Kubernetes cluster with LLM-generated
input criteria.

\textbf{Positioning.} Table~\ref{tab:related-comparison} compares Ai4k8s
against the most relevant systems.

\begin{table}[h]
  \centering
  \caption{Comparison of autoscaling and AIOps systems.}
  \label{tab:related-comparison}
  \small
  \begin{tabular}{lcccccc}
    \toprule
    System & Predictive & LLM & MCDA & Explainable & On-premises & Open-source \\
    \midrule
    HPA (Kubernetes) & \texttimes & \texttimes & \texttimes & \texttimes & \checkmark & \checkmark \\
    VPA (Kubernetes) & \texttimes & \texttimes & \texttimes & \texttimes & \checkmark & \checkmark \\
    K8sGPT          & \texttimes & \checkmark & \texttimes & \checkmark & partial & \checkmark \\
    AIOpsLab        & \texttimes & \checkmark & \texttimes & \checkmark & \texttimes & \checkmark \\
    \textbf{Ai4k8s (ours)} & \checkmark & \checkmark & \checkmark & \checkmark & \checkmark & \checkmark \\
    \bottomrule
  \end{tabular}
\end{table}
